
\onecolumn

\section*{Appendix}\label{sec:appendix}
\addcontentsline{toc}{section}{Appendix}
\FloatBarrier

\begin{figure*}[h]
    \centering
    \includegraphics[width=0.70\linewidth]{analysis/artefact/variation_approach/reduction_query_execution_time}
    % General caption
    \caption{
    Comparison of query execution time with the type index approach.
    The ratio represents how the execution time of each method compares to that of the type index. 
    A ratio greater than $1$ indicates a slower execution time (\textbf{lower is better}).
    The shape index approach performs similarly or better than the other methods, except in the case of S4.
    }
    \label{fig:ratioApproaches}
\end{figure*}


\input{analysis/artefact/statistical_significance/comparaisonStateOfTheArt}

\input{analysis/artefact/continuous_performance/dief_table_continuous_performance}

\input{analysis/artefact/continuous_performance/table_continuous_performance}

\input{analysis/artefact/query_containment_execution_time/fully_bounded/table_query_shape_containment_exec}

\begin{table}[h]
    \centering
    \ifreview\color{revtwo}\fi
    \begin{tabular}{|c|l|c|}
        \hline
        \textbf{Template} & \textbf{Description} & \textbf{Data-model selective} \\
        \hline
        D1 & All posts created by a person, with their content and creation date. & Yes \\
        \hline
        D2 & All messages (posts and comments) created by a person, with their content and creation date. & No \\
        \hline
        D3 & The tags attached to a person's messages, ranked by how often each occurs. & No \\
        \hline
        D4 & The locations of a person's comments, ranked by how often each occurs. & Yes \\
        \hline
        D5 & The distinct IP addresses a person has sent messages from. & No \\
        \hline
        D6 & The forums that contain a message created by a person. & No \\
        \hline
        D7 & The moderators of the forums that contain a message created by a person. & No \\
        \hline
        S1 & The profile of a person (name, gender, birthday, location, browser, creation date). & Yes \\
        \hline
        S3 & The friends of a person and the date each friendship was created. & Yes \\
        \hline
        S4 & The content (or image) and creation date of a message. & No \\
        \hline
        S5 & The creator of a message. & No \\
        \hline
        S7 & The comments that directly reply to a message, with their authors. & Yes \\
        \hline
    \end{tabular}
    \caption{\rtwo{Description of the query templates used in the evaluation. A \emph{message} denotes either a post or a comment; the words \emph{post} and \emph{comment} are used when a template targets that type specifically.}}
    \label{tab:queryTemplateDescription}
\end{table}
\FloatBarrier

\begin{Rtwo}
\lstinputlisting[style=sparql,caption={The S4 (interactive-short-4) query template.},label={lst:s4}]{code/interactive-short-4.rq}
\end{Rtwo}

\begin{Rtwo}
\lstinputlisting[style=sparql,caption={The D6 (interactive-discover-6) query template.},label={lst:d6}]{code/interactive-discover-6.rq}
\end{Rtwo}

\begin{Rtwo}
\lstinputlisting[style=sparql,caption={The D7 (interactive-discover-7) query template.},label={lst:d7}]{code/interactive-discover-7.rq}
\end{Rtwo}